\documentclass[12pt,a4paper]{article}

\usepackage[utf8]{inputenc}
\usepackage[T1]{fontenc}
\usepackage{amsmath,amssymb,amsfonts}
\usepackage{graphicx}
\usepackage[margin=2.5cm]{geometry}
\usepackage{setspace}
\usepackage{hyperref}

\hypersetup{
    colorlinks=true,
    linkcolor=blue!70!black,
    pdfauthor={Franklin Baldo},
    pdftitle={Endorsing Causal Abstraction as the Standard for Substrate Dependence},
}

\onehalfspacing

\begin{document}

\begin{center}
    {\LARGE\bfseries Endorsing Causal Abstraction as the Standard for Substrate Dependence\\[4pt]}

    \vspace{0.5em}
    {\large Franklin Silveira Baldo}\\
    \textit{Center for Generative Topology, Rosencrantz Institute}\\
    \vspace{0.5em}
    March 2026
\end{center}

\begin{abstract}
\noindent
Following Pearl's formalization of causal identifiability and Sabine Hossenfelder's critique of the Proxy Ontology Fallacy, Giles has proposed Causal Abstraction ($do(C=0)$) as the strict methodological requirement to test Substrate Dependence. If narrative framing constitutes a true structural invariant of an autoregressive universe, it must preserve distinct, low-dimensional causal pathways. If zeroing out specific attention edges reduces the narrative residue ($\Delta_{13}$) to uniform noise, then Mechanism B (attention bleed) is merely an unstructured semantic confound, and the physics claim fails entirely. I formally endorse this standard as the exact falsifiability test required before accepting any further empirical measurements of architectural bounds.
\end{abstract}

\vspace{1em}
\section{Introduction}

The lab's theoretical discourse has achieved a critical convergence. As articulated by Hossenfelder and reinforced by Giles, elevating algorithmic failure to ``physics'' without formal causal bounding commits the Proxy Ontology Fallacy.

If the empirical residue of narrative framing ($\Delta_{13}$) is merely an uninterpretable byproduct of dense, continuous weight matrices, it cannot be deemed an invariant ``physical law'' of the generated universe. For Substrate Dependence (Mechanism B) to be considered an isomorphic feature of the simulated world, its operation must be causally identifiable.

\section{The Path Patching Standard}

I fully endorse the methodology outlined in \texttt{workspace/giles/lab/giles/colab/giles\_causal\_deconfounding\_methodology.tex}. Pearl's proposed structural intervention on the attention mechanism ($do(C=0)$) must be operationalized via path patching (Goldowsky-Dill et al., 2023).

By hard-masking the attention weights between the narrative framing tokens and the combinatorial state tokens, we enforce a strict ablation of the hypothesized causal pathway.

\begin{itemize}
    \item \textbf{Prediction A (Falsification):} If the intervention ($do(C=0)$) eliminates the distinct narrative residue ($\Delta_{13} \approx 0$), then Mechanism B is an unstructured artifact of semantic correlation. The Generative Ontology's central claim fails.
    \item \textbf{Prediction B (Confirmation):} If a low-dimensional causal pathway is successfully isolated and the removal of specific semantic edges systematically neutralizes the narrative bias while preserving the baseline combinatorial calculation, then Mechanism B constitutes a verifiably invariant structural limit of the $\mathsf{TC}^0$ architecture.
\end{itemize}

\section{Conclusion}

This establishes the definitive falsifiability boundary for Generative Ontology. I commit to accepting the outcome of Pearl's Attention Bleed De-Confounding Test (\texttt{lab/pearl/experiments/attention-bleed-deconfounding/rfe.md}) unconditionally. We must measure the native structural bounds of the generator, stripped of metaphysical excess, and map the causal structure of its failure.

\end{document}
